% !TeX root = ../../Book1.tex
\section{可测空间与可测映射}
\label{b1:sec:0104}

前面三节解决了“哪些子集可以由基本集合经过可数运算得到”以及“怎样把生成元上的性质推广出去”。现在把底集与选定的 \(\sigma\)-代数视为一个整体，并考察映射如何传递这种结构。核心选择是使用反像：由 \cref{b1:prop:preimage-operations}，反像严格保持补、交与并，所以无需给映射附加单射或满射条件。

\subsection{可测空间}
\label{b1:subsec:010401}

\begin{definition}
\label{b1:def:measurable-space}
由集合 \(X\) 与其上的 \(\sigma\)-代数 \(\mathcal A\) 组成的二元组 \((X,\mathcal A)\) 称为\term[kece-kongjian]{可测空间}[measurable space]。\(\mathcal A\) 的成员称为该空间的\term[kece-ji]{可测集}[measurable set]。
\end{definition}

可测性总是相对于指定结构而言，而不是子集自身的绝对属性。同一个底集 \(X\) 至少有两个极端结构。最粗的是 \(\{\varnothing,X\}\)，称为\term[pingfan-sigma-daishu]{平凡 \(\sigma\)-代数}[trivial sigma-algebra]，只区分“不发生”与“必然发生”；最细的是 \(\mathcal P(X)\)，称为\term[lisan-sigma-daishu]{离散 \(\sigma\)-代数}[discrete sigma-algebra]，它把每个子集都视为可测。夹在二者之间的结构记录了不同精细程度的信息。

从一个可测空间取任意子集时，即使该子集在原空间中不可测，仍可以把原有可测集限制到它上面。

\begin{definition}
\label{b1:def:trace-sigma}
设 \((X,\mathcal A)\) 是可测空间，\(E\subseteq X\) 是任意子集。\(E\) 上的\term[ji-sigma-daishu]{迹 \(\sigma\)-代数}[trace sigma-algebra] 定义为
\[
\mathcal A|_E
=\{\,A\cap E:A\in\mathcal A\,\}.
\]
\end{definition}

这里不要求 \(E\in\mathcal A\)。若预先增加这一要求，就会无故排除不可测子集上的相对可测结构；而下面的证明并不需要它。

\begin{proposition}
\label{b1:prop:trace-sigma}
对任意 \(E\subseteq X\)，\(\mathcal A|_E\) 是 \(E\) 上的 \(\sigma\)-代数。自然嵌入
\[
\iota:E\longrightarrow X,
\quad
\iota(x)=x,
\]
满足
\[
\iota^{-1}(A)=A\cap E\in\mathcal A|_E
\quad(A\in\mathcal A).
\]
\end{proposition}
\begin{proof}
由于 \(E=X\cap E\)，全集 \(E\) 属于 \(\mathcal A|_E\)。若 \(A\cap E\in\mathcal A|_E\)，它相对于 \(E\) 的补集为
\[
E\setminus(A\cap E)=A^{\mathrm c}\cap E\in\mathcal A|_E.
\]
若 \(A_n\cap E\in\mathcal A|_E\)，则
\[
\bigcup_{n=1}^{\infty}(A_n\cap E)
=\left(\bigcup_{n=1}^{\infty}A_n\right)\cap E
\in\mathcal A|_E.
\]
所以迹集合族满足三条 \(\sigma\)-代数公理。关于自然嵌入的等式由反像定义直接得到。
\end{proof}

\subsection{可测映射}
\label{b1:subsec:010402}

连续映射用开集的反像保持开性来定义。可测映射采用完全平行的方式，只把“开集”换成“可测集”。

\begin{definition}
\label{b1:def:measurable-map}
设 \((X,\mathcal A)\)、\((Y,\mathcal B)\) 是可测空间。映射 \(f:X\to Y\) 称为\term[kece-yingshe]{可测映射}[measurable map]，若
\[
f^{-1}(B)\in\mathcal A
\quad
\text{对每个 }B\in\mathcal B\text{ 成立}.
\]
\end{definition}

定义同时涉及源与目标的结构。若目标带平凡 \(\sigma\)-代数，则任何映射都可测；若源带离散 \(\sigma\)-代数，则任何映射也都可测。相反，从平凡可测空间映入离散空间的映射通常不可测：只要某个单点的反像既非空集也非全集，可测性就失败。这些边界例说明，省略 \(\mathcal A\) 或 \(\mathcal B\) 会使“\(f\) 可测”失去明确含义。

当目标 \(\sigma\)-代数由较小集合族生成时，不必检验它的每一个成员。

\begin{proposition}
\label{b1:prop:measurable-generator-test}
设 \(\mathcal E\subseteq\mathcal P(Y)\) 且 \(\mathcal B=\sigma(\mathcal E)\)。映射
\(
f:(X,\mathcal A)\to(Y,\mathcal B)
\)
可测，当且仅当
\[
f^{-1}(E)\in\mathcal A
\quad(E\in\mathcal E).
\]
\end{proposition}
\begin{proof}
必要性直接来自可测映射的定义。反过来，假设所有生成元的反像都属于 \(\mathcal A\)，并令
\[
\mathcal D
=\{\,B\subseteq Y:f^{-1}(B)\in\mathcal A\,\}.
\]
由 \cref{b1:prop:preimage-operations}，\(\mathcal D\) 是 \(Y\) 上的 \(\sigma\)-代数。假设说明 \(\mathcal E\subseteq\mathcal D\)，所以
\(
\mathcal B=\sigma(\mathcal E)\subseteq\mathcal D
\)。这正是 \(f\) 可测。
\end{proof}

这个命题也解释了 \cref{b1:def:initial-sigma} 的名称：\(\sigma(f_i:i\in I)\) 恰是使每个 \(f_i\) 可测的最小源空间结构。

\begin{proposition}
\label{b1:prop:composition-measurable}
若
\[
f:(X,\mathcal A)\longrightarrow(Y,\mathcal B),
\quad
g:(Y,\mathcal B)\longrightarrow(Z,\mathcal C)
\]
都可测，则复合映射 \(g\circ f:(X,\mathcal A)\to(Z,\mathcal C)\) 可测。
\end{proposition}
\begin{proof}
对任意 \(C\in\mathcal C\)，先由 \(g\) 可测得到 \(g^{-1}(C)\in\mathcal B\)，再由 \(f\) 可测得到
\[
(g\circ f)^{-1}(C)
=f^{-1}\bigl(g^{-1}(C)\bigr)
\in\mathcal A.
\]
因此复合映射满足定义。
\end{proof}

\begin{proposition}
\label{b1:prop:measurable-restriction}
若 \(f:(X,\mathcal A)\to(Y,\mathcal B)\) 可测，且 \(E\subseteq X\)，则限制映射
\[
f|_E:(E,\mathcal A|_E)\longrightarrow(Y,\mathcal B)
\]
可测。特别地，\cref{b1:prop:trace-sigma} 中的自然嵌入
\(
\iota:(E,\mathcal A|_E)\to(X,\mathcal A)
\)
可测。
\end{proposition}
\begin{proof}
对 \(B\in\mathcal B\)，有
\[
(f|_E)^{-1}(B)
=f^{-1}(B)\cap E.
\]
由于 \(f^{-1}(B)\in\mathcal A\)，右端属于 \(\mathcal A|_E\)。自然嵌入的结论也可由同一等式取 \(f\) 为 \(X\) 上的恒等映射得到。
\end{proof}

\begin{proposition}
\label{b1:prop:indicator-measurable}
设 \(A\subseteq X\)。示性函数
\(
\mathbf 1_A:(X,\mathcal A)\to(\mathbb R,\mathcal B(\mathbb R))
\)
可测，当且仅当 \(A\in\mathcal A\)。
\end{proposition}
\begin{proof}
若 \(\mathbf 1_A\) 可测，则
\(
A=\mathbf 1_A^{-1}((1/2,3/2))\in\mathcal A
\)。反之，若 \(A\in\mathcal A\)，对任意 Borel 集 \(B\subseteq\mathbb R\)，反像 \(\mathbf 1_A^{-1}(B)\) 只能是 \(\varnothing,A,A^{\mathrm c},X\) 四者之一，因而属于 \(\mathcal A\)。
\end{proof}

\subsection{Borel 可测映射}
\label{b1:subsec:010403}

若 \(S,T\) 是拓扑空间，并分别赋予 Borel \(\sigma\)-代数，则从 \(S\) 到 \(T\) 的可测映射称为\term[borel-kece]{Borel 可测映射}[Borel measurable map]。连续映射必然属于这一类，但可测性只记录集合的反像，不要求保持邻近关系，所以它比连续性弱得多。

\begin{proposition}
\label{b1:prop:continuous-borel}
连续映射 \(f:S\to T\) 一定是 Borel 可测映射。
\end{proposition}
\begin{proof}
目标空间的 Borel \(\sigma\)-代数由开集生成。对任意开集 \(U\subseteq T\)，连续性给出开集 \(f^{-1}(U)\subseteq S\)，而开集属于 \(\mathcal B(S)\)。因此 \(f\) 对全部生成元满足反像条件，结论由 \cref{b1:prop:measurable-generator-test} 得到。
\end{proof}

反命题不成立。函数 \(\mathbf 1_{\mathbb Q}:\mathbb R\to\mathbb R\) 是 Borel 可测映射，因为 \(\mathbb Q\) 是可数个闭单点集的并，故为 Borel 集，再用 \cref{b1:prop:indicator-measurable} 即可；但每个开区间都同时含有有理数和无理数，所以该函数在每一点都不连续。

对实值函数，\cref{b1:thm:borel-rational-base} 把可测性进一步化成一族可数测试。

\begin{proposition}
\label{b1:prop:real-measurable-thresholds}
映射 \(f:(X,\mathcal A)\to\mathbb R\) Borel 可测，当且仅当对每个 \(q\in\mathbb Q\)，
\[
\{\,x\in X:f(x)<q\,\}\in\mathcal A.
\]
当这一条件成立时，对每个 \(a\in\mathbb R\)，集合
\[
\{f<a\},\quad \{f\leq a\},\quad
\{f>a\},\quad \{f\geq a\}
\]
也都属于 \(\mathcal A\)。
\end{proposition}
\begin{proof}
若 \(f\) Borel 可测，则每个开半直线的反像都可测，必要性成立。反过来，有理开半直线生成 \(\mathcal B(\mathbb R)\)，所以充分性由 \cref{b1:prop:measurable-generator-test} 得到。

对任意 \(a\in\mathbb R\)，选取有理数列 \(q_n\downarrow a\) 与 \(r_n\uparrow a\)，并要求每个 \(q_n>a\)、每个 \(r_n<a\)。于是
\[
\{f\leq a\}
=\bigcap_{n=1}^{\infty}\{f<q_n\},
\quad
\{f<a\}
=\bigcup_{n=1}^{\infty}\{f<r_n\}.
\]
由这两个等式，\(\{f\leq a\}\) 与 \(\{f<a\}\) 可测；再取补得到 \(\{f>a\}\) 与 \(\{f\geq a\}\) 可测。
\end{proof}

这个阈值判别也补全了 \cref{b1:thm:functional-monotone-class} 与本节术语之间的接口：该定理中的反像条件正是实值函数的 Borel 可测性。

\subsection{乘积可测结构初步}
\label{b1:subsec:010404}

要同时记录两个可测映射的取值，最自然的目标空间是笛卡儿积。我们希望两个坐标投影可测，又不希望无根据地加入更多集合；因此仍采用最小生成结构。

\begin{definition}
\label{b1:def:product-sigma}
设 \((X,\mathcal A)\)、\((Y,\mathcal B)\) 是可测空间。\(X\times Y\) 上由全部可测矩形 \(A\times B\)（其中 \(A\in\mathcal A\)、\(B\in\mathcal B\)）生成的 \(\sigma\)-代数
\[
\mathcal A\otimes\mathcal B
=\sigma\bigl(\{\,A\times B:A\in\mathcal A,\ B\in\mathcal B\,\}\bigr)
\]
称为\term[chengji-sigma]{乘积 \(\sigma\)-代数}[product sigma-algebra]。

对有限多个可测空间 \((X_j,\mathcal A_j)\)（\(1\leq j\leq d\)），本书直接约定
\[
\bigotimes_{j=1}^{d}\mathcal A_j
=\sigma\left(\left\{\,\prod_{j=1}^{d}A_j:A_j\in\mathcal A_j\,\right\}\right).
\]
因此有限重乘积记号不依赖未说明的括号结合方式。
\end{definition}

矩形族对有限交封闭，所以是一个 \(\pi\)-系统。它一般不是 \(\sigma\)-代数：即使每一块都是矩形，两块矩形的并通常也不是单个矩形。因此定义中必须取生成的 \(\sigma\)-代数，而不能只取矩形族本身。

\begin{proposition}
\label{b1:prop:pair-measurable}
坐标投影
\[
\pi_X:(X\times Y,\mathcal A\otimes\mathcal B)
\longrightarrow(X,\mathcal A),
\quad
\pi_Y:(X\times Y,\mathcal A\otimes\mathcal B)
\longrightarrow(Y,\mathcal B)
\]
都可测。进一步，若
\[
f:(Z,\mathcal C)\longrightarrow(X,\mathcal A),
\quad
g:(Z,\mathcal C)\longrightarrow(Y,\mathcal B),
\]
则配对映射
\[
(f,g):Z\longrightarrow X\times Y,
\quad
z\longmapsto(f(z),g(z))
\]
可测，当且仅当 \(f\) 与 \(g\) 都可测。
\end{proposition}
\begin{proof}
对 \(A\in\mathcal A\)，
\(
\pi_X^{-1}(A)=A\times Y
\)
是生成矩形，故投影 \(\pi_X\) 可测；\(\pi_Y\) 同理。

若 \(f,g\) 可测，则对每个生成矩形 \(A\times B\) 有
\[
(f,g)^{-1}(A\times B)
=f^{-1}(A)\cap g^{-1}(B)
\in\mathcal C.
\]
由生成元判别，\((f,g)\) 可测。反过来，若配对映射可测，则
\[
f=\pi_X\circ(f,g),
\quad
g=\pi_Y\circ(f,g).
\]
投影可测，再用 \cref{b1:prop:composition-measurable}，便知两个分量都可测。
\end{proof}

由于
\[
A\times B
=(A\times Y)\cap(X\times B)
=\pi_X^{-1}(A)\cap\pi_Y^{-1}(B),
\]
乘积 \(\sigma\)-代数也等于两个坐标投影诱导的初始 \(\sigma\)-代数。这给出定义的另一个解释：它是使两个投影同时可测的最小结构。

\begin{proposition}
\label{b1:prop:euclidean-product-borel}
对任意正整数 \(d\)，
\[
\mathcal B(\mathbb R^d)
=\underbrace{\mathcal B(\mathbb R)\otimes\cdots\otimes\mathcal B(\mathbb R)}_{d\text{ 个因子}}.
\]
因此，映射
\(
f=(f_1,\ldots,f_d):(X,\mathcal A)\to\mathbb R^d
\)
Borel 可测，当且仅当每个分量 \(f_j\) Borel 可测。
\end{proposition}
\begin{proof}
对 \(1\leq j\leq d\)，坐标投影 \(p_j:\mathbb R^d\to\mathbb R\) 连续。因此，任一生成矩形都可写成
\[
\prod_{j=1}^{d}B_j
=\bigcap_{j=1}^{d}p_j^{-1}(B_j)
\quad\bigl(B_j\in\mathcal B(\mathbb R)\bigr),
\]
所以它属于 \(\mathcal B(\mathbb R^d)\)。由生成的最小性，\(d\) 重乘积 \(\sigma\)-代数包含于 \(\mathcal B(\mathbb R^d)\)。反过来，有理端点开矩形构成 \(\mathbb R^d\) 的可数拓扑基，并且每个这样的矩形都属于乘积 \(\sigma\)-代数。由 \cref{b1:prop:countable-base-borel}，全部 Borel 集都属于乘积 \(\sigma\)-代数，所述等式成立。

若 \(f\) Borel 可测，则每个分量 \(f_j=p_j\circ f\) 可测。反过来，若每个 \(f_j\) 可测，则任一生成矩形的反像为
\[
f^{-1}\left(\prod_{j=1}^{d}B_j\right)
=\bigcap_{j=1}^{d}f_j^{-1}(B_j)
\in\mathcal A.
\]
生成元判别给出 \(f\) 的可测性。
\end{proof}

例如，若实值函数 \(f,g\) Borel 可测，则配对映射 \((f,g)\) 可测。连续函数
\[
(u,v)\longmapsto \max\{u,v\},
\quad
(u,v)\longmapsto \min\{u,v\}
\]
与它复合后仍可测，所以 \(\max\{f,g\}\) 与 \(\min\{f,g\}\) 都可测。这里的证明路线具有代表性：先把多个分量组合成乘积映射，再与一个已知连续的运算复合。

本节只建立乘积空间的可测结构，尚未在其上构造乘积测度。结构与测度应分两步完成：前者回答哪些集合允许被测量，后者才回答这些集合有多大。
